\section{Theoretical Background}

This section establishes the conceptual foundations on which the review rests,
covering the contemporary challenges of supply chain management, the main
risk categories recurring across the literature, and the AI techniques most
relevant to the studies examined.
 
\subsection{Supply Chain Management: Scope and Contemporary Challenges}
 
Supply Chain Management (SCM) encompasses the coordination of activities
involved in sourcing raw materials, transforming them into finished goods, and
delivering them to end customers, together with the management of material,
information, and financial flows that support these processes. For much of the
2000s, lean management was the dominant paradigm: reducing inventory
buffers, rationalising supplier bases, and optimising for cost. These principles
delivered genuine efficiency gains but created systemic fragility that went
largely unnoticed until disruptions struck at scale.
 
The COVID-19 pandemic exposed this fragility with particular severity.
Manufacturing firms dependent on single-source strategies found themselves
unable to absorb even brief supply shocks, with shortages persisting well into
2021. Similar dynamics had emerged after the 2011 Tohoku earthquake and
the 2021 Suez Canal blockage. These repeated crises accelerated a conceptual
shift already visible in the academic literature, with researchers arguing that
supply chains need to be designed not merely for efficiency but for
\emph{resilience}, the capacity to anticipate disruptions, absorb their
impact, and recover within a reasonable timeframe, and, in the case of
prolonged crises, for \emph{viability}, meaning the capacity to adapt
structurally rather than simply return to a prior steady state.
 
The globalisation of supply chains has also introduced a complexity that
compounds these vulnerabilities. Modern value networks typically span
multiple supplier tiers, logistics intermediaries, and regulatory environments
across continents. Interdependencies between nodes mean that a disruption at
one point can propagate rapidly through the network, the so-called
\emph{ripple effect}, making containment one of the central problems to
which AI-based approaches are being applied.
 
\subsection{Supply Chain Risk: Categories and Taxonomy}
 
Supply chain risk is broadly understood as the potential for an event to
disrupt material, information, or financial flows in ways that harm
organisational performance. This review adopts a multi-category taxonomy
drawing on the typologies most frequently appearing across the included
studies.
 
\textbf{Supplier risk} refers to the possibility that a supplier will fail to
deliver inputs on time, at the required quantity, or quality level, and
encompasses financial instability, capacity constraints, and geographic
concentration. \textbf{Demand risk} arises from uncertainty in customer
orders, including the well-documented bullwhip effect whereby small retail
fluctuations are amplified into large upstream production swings.
\textbf{Disruption risk} covers low-frequency, high-impact external events,
pandemics, natural disasters, geopolitical conflicts, port closures, that
are poorly represented in historical data. \textbf{Operational risk}
encompasses internally originating failures such as equipment breakdowns,
process errors, and IT system failures. \textbf{Financial and fraud risk},
including payment defaults and invoice manipulation, receives comparatively
less attention in the corpus but is identified as a clear use case for
explainable AI. Two further categories, environmental and sustainability
risk and cybersecurity risk, are of growing prominence in the broader
literature, though they remain underrepresented in this corpus.
 
\subsection{Artificial Intelligence in Supply Chain Management}
 
Artificial intelligence, in the context of this review, denotes computational
methods that enable systems to learn from data, detect patterns, and support
or automate decisions. The corpus reflects a broad methodological spectrum.
 
\textbf{Traditional machine learning}, including Random Forest, Support
Vector Machines, and gradient boosting methods such as XGBoost and
LightGBM, remains the most prevalent approach, valued for its robustness on
structured tabular data and its interpretability relative to deeper
architectures. Unsupervised methods, particularly K-Means clustering, are
applied to supplier segmentation and anomaly detection.
 
\textbf{Deep learning} architectures, foremost Long Short-Term Memory
networks (LSTM) and Gated Recurrent Units (GRU), are applied to demand
forecasting, disruption prediction, and time-series anomaly detection.
Hybrid architectures combining LSTMs with Graph Neural Networks (GNNs)
model both the temporal evolution of risk indices and the structural
relationships between supply chain nodes. One notable study develops a
sequence-to-disruption architecture integrating bidirectional GRUs with
transformer-style attention, achieving prediction errors under 20~days for
six-month-ahead shortage forecasts on a dataset of 10,000 automotive
components.
 
\textbf{Graph Neural Networks} operate on graph-structured supply chain
data, learning how disruptions propagate through network topology. One
reviewed framework constructs a supply chain knowledge graph with over
2.5~million entities and eleven relationship types to support real-time risk
monitoring and recovery decisions.
 
\textbf{Natural language processing} enables structured information to be
extracted from unstructured text. Applications include sentiment analysis of
social media streams for real-time early warning and named entity recognition
pipelines applied to longitudinal news archives for risk categorisation.
 
\textbf{Explainable AI (XAI)}, particularly SHAP (SHapley Additive
Explanations), addresses the interpretability deficit of complex models. Its
adoption reflects a practical imperative: supply chain risk decisions carry
significant operational and financial consequences, requiring not only
accurate predictions but transparent justifications that managers can
interrogate and audit.
 
\textbf{Federated machine learning} addresses the competitive and regulatory
barriers to centralising sensitive supplier performance data by allowing
organisations to collaboratively train shared risk models without raw data
leaving individual servers. Studies in the corpus indicate that federated
models achieve performance close to centralised counterparts while preserving
data privacy, a trade-off particularly attractive for small and medium-sized
enterprises.